\chapter[How to Tune]{Tuning your strings}

Once you have all your supplies prepared and understand how to care for them, there remains one last task before we can begin "playing" the guitar. You must learn how to tune the strings.

\section{What exactly is tuning?}

Tuning the strings refers to tightening or loosening the strings of the guitar with the tuning keys to make each one produce a specific pitch when plucked. An electronic tuner or tuning app makes this task much less daunting than it sounds.

\section{Why do I need to tune?}

Tuning is what makes the notes of the guitar sound ``correct.'' If you do not tune your guitar correctly, everything you play will sound strange and wrong. Also, you will not be able to play with other people until you learn how to tune.

\section[Proper tuning]{What all do I need to know in order to tune properly?}

The challenge of tuning comes from the complexity of the task - there are several different things you must keep track of while tuning.

\subsection{The tuning keys}

Each string has a corresponding tuning key on the headstock of the guitar that controls its tension, or how tight or loose the string is. Changing tension is exactly what we are doing when we are tuning our strings. However, all the keys look identical, so you must take care to memorize which key goes with which string. Figure 2.1 shows which key goes with which string. Pay attention to how each string looks slightly different. The nylon strings are clear, the metal strings are obviously metal-wrapped, and each string has a different thickness.

\begin{figure}
	
	\centering
	\includegraphics[width=1\linewidth]{../../Downloads/headstock}
	\caption{Notice the order of the strings}
	\label{fig:headstock}
\end{figure}

\subsection{Tight vs. loose and sharp vs. flat}

Each tuning key controls the \emph{tension} of the string, which is the feeling of how tight or loose the string is. Different tensions produce different notes. In music, however, we prefer to use the words \emph{sharp} and \emph{flat} rather than tight or loose.

Your tuner may show you that your string is sharp, which means too tight. You will need to loosen the string accordingly to tune it. A too-loose string is called flat and tightened to the correct pitch.  

\subsection{How do I know which direction to turn?}

The tuning key controls the rotation of a small cylinder in the headstock that the string is attached to. Gently turn the key one direction or the other and notice the direction of rotation of the cylinder. If it rotates clockwise towards the bridge, the string is being loosened and will get flatter. If it rotates counter-clockwise towards the top of the headstock, the string is being stretched tighter and will get sharper.

\subsection{What is my tuner telling me?}

Your tuner can tell one of three things, that the note played is either in tune, too sharp, or too flat. In tune notes are typically indicated by a green or blue color, and will often show a color or icon centered in the screen to indicate that the note is in tune. Most electronic tuners and tuning apps show sharp vs. flat on either a horizontal or vertical spectrum. 

\vspace{5mm}

{\large Flat} {\Huge E} {\large Sharp}







\section{How do I tune my strings?}

Tuning is specific set of steps you must repeat on every string. Eventually, this will become automatic as you do it every day.

\begin{itemize}
	\item Prepare your electronic tuner or tuning app. Most electronic tuners need to be clipped on to the guitar headstock. Most tuning apps use your phone's microphone to listen to the guitar, so place your phone close to the sound hole.
	\item Play the first string
\end{itemize}


